Electric vehicles are spreading fast, and the grid has noticed. Public
charging stations behave like a new layer of distributed load that is
neither uniform nor predictable: a single station can swing from idle
on a Saturday to peak demand on a Tuesday morning, and the swings carry
weekly seasonality, weather sensitivity and the occasional holiday
shock. For grid operators and station owners this matters because
weekly capacity planning, maintenance scheduling and energy procurement
all benefit from forecasting demand at least a week ahead. That is the
problem we picked.

Boulder, Colorado is a convenient laboratory for this question. The
city has run a network of 26 public level-2 stations since 2018, and
it publishes every charging session through its Open Data portal.
The dataset we work with contains 148\,136 individual sessions logged
between August 2018 and November 2023 --- five and a half years of
real-world data, which is long enough to capture multiple annual
cycles, the COVID-19 disruption and the post-pandemic recovery. Each
record carries the station name, the start and end timestamps, the
charging duration, the energy delivered (kWh) and the estimated
greenhouse gas savings. We aggregate by station-day, focus on the
station with the highest session volume (N~Boulder Rec~1, 16\,278
sessions), and treat the resulting daily kWh series as our forecasting
target.

The horizon is seven days ahead. We chose a week because it covers
a full weekday-weekend cycle and because most distribution utilities
plan on a rolling-week basis, so a forecast on this scale is what an
operator would actually use. Our initial scope was the three families
required by the course: naive baselines, the (S)ARIMA family, and an
LSTM. As we worked through them we kept going: we ran a head-to-head
between two zero-shot foundation models (Google's TimesFM and Amazon's
Chronos), and finally implemented a PatchTST transformer from scratch
in PyTorch with a three-seed ensemble. We did not implement
CAT-Former, the closest published state of the art on this dataset,
but we read the paper and take it as a benchmark for the discussion.

This paper documents the dataset, the protocol, every model we ran
under that protocol, and the claim we can defensibly make about
beating the state of the art. We beat the strongest model in
our own leaderboard. We do not claim to beat the published numbers,
which use different units and station selections --- the full
discussion is in Sections~\ref{sec:results}
and~\ref{sec:conclusions}.

The paper is organised as follows.
Section~\ref{sec:related} reviews the state of the art on EV
charging demand forecasting.
Section~\ref{sec:data} describes the dataset and its statistical
properties.
Section~\ref{sec:method} defines the train/test protocol and the
metrics.
Section~\ref{sec:models} walks through every model family we ran.
Section~\ref{sec:results} gives the full leaderboard, the per-window
analysis, and the SoTA-claim breakdown.
Section~\ref{sec:conclusions} closes with conclusions, lessons learned
and limitations.

\medskip
\noindent
\textbf{What is our contribution?} The PatchTST architecture itself is
not ours --- it comes from Nie et al.~\cite{nie2023patchtst}. What we
did add is: (i)~an implementation from scratch in PyTorch that lives
in the project repository, with RevIN~\cite{kim2022revin} on the
target channel and channel-mixed patching over four input channels
(energy, temperature, weekend flag, holiday flag); (ii)~a three-seed
ensemble strategy that cuts inference variance by about~7\% MAE; and
(iii)~a simple-mean ensemble with the LSTM that we picked
\emph{before} looking at the test scores, which gave the best model
overall. All code, data and trained-model artefacts are openly
available.\footnote{Project repository:
\href{https://github.com/eeguskiza/ev-charging-forecasting}{\texttt{github.com/eeguskiza/ev-charging-forecasting}}~\cite{ourrepo}}
